\section{Навье--Стокс: каскад, интеграл и диадическая модель}\label{sec:ns}

Навье--Стокс входит в статью по одной структурной причине, и оговорку мы делаем сразу и громко:
\emph{никакой связи с простыми числами здесь нет и не заявляется}. Ветка не кормит близнецов и не
кормится от них. Общее у неё с ними --- единственный лозунг: дефект не исчезает, и если он не
закрылся, за него надо платить, --- реализованный тут в самой буквальной валюте, вязкой диссипации
энергии. На этом настоящем аналитическом объекте мы показываем, докуда честно достаёт бюджетная
дисциплина и где нетронутой стоит проблема тысячелетия. Уравнение формализовано в сильной форме на
$E_3 := \code{EuclideanSpace}\ \R\ (\code{Fin }3)$, а дифференциальные операторы (дивергенция,
векторный лапласиан, конвективный член) собраны из голого \code{fderiv}.

\begin{definition}[\code{IsNSSolution}]\label{def:IsNSSolution}
Для вязкости $\nu$, внешней силы $f$, поля скоростей $u$ и давления $p$
\begin{equation}\label{eq:ns-solution}
\code{IsNSSolution}\ \nu\ f\ u\ p \;:=\;
\bigl[\partial_t u + (u\cdot\nabla)u = \nu\,\Delta u - \nabla p + f\bigr]
\;\wedge\; \bigl[\operatorname{div} u = 0\bigr].
\end{equation}
\end{definition}

\noindent Это несжимаемые уравнения Навье--Стокса \cite{clay-ns,millennium} в классической сильной
форме. Никаких слабых решений и распределений: сильная форма выбрана сознательно, чтобы предикат
читался глазами. Первый вопрос к любому такому предикату --- не пуст ли он; ответ машинный.

\begin{theorem}[\code{zero\_is\_NSSolution}]\label{thm:zero_is_NSSolution}
$\code{IsNSSolution}\ \nu\ 0\ 0\ 0$ для любого $\nu \in \R$: нулевое поле с нулевым давлением и без
силы есть решение при любой вязкости. \green{}
\end{theorem}

\emph{Идея доказательства.} Прямое вычисление; каждый член \cref{eq:ns-solution} обращается в нуль.
Скромно, но это ровно прививка от вакуумности: предикат обитаем, уравнение настоящее.

Энергетика задаётся интегралами Бохнера по объёму: кинетическая энергия
$\code{kineticEnergy} = \tfrac12\int\|u\|^2$ и скорость диссипации
$\code{dissipationRate} = \nu\int\sum_i\|\partial_i u\|^2$ (энстрофийная форма). Здесь спрятана
ловушка, и репозиторий вскрывает её машинно: интеграл Бохнера в \code{mathlib} \emph{молча равен
нулю} на неинтегрируемой функции (\code{MeasureTheory.integral\_undef}).

\begin{theorem}[\code{kineticEnergy\_of\_not\_integrable}]\label{thm:kineticEnergy_of_not_integrable}
Если $\|u\|^2$ не интегрируемо (нет предиката \code{FiniteKineticEnergy}), то
$\code{kineticEnergy}\ u = 0$. \green{}
\end{theorem}

\emph{Идея доказательства.} Переписывание через \code{MeasureTheory.integral\_undef}. Равенство
$\code{kineticEnergy}\ u = 0$ тем самым может означать не «энергия нулевая», а «энергия бесконечна
или не определена»; «нулевая энергия» неинтегрируемого поля --- артефакт определения, не физика.
Мораль: всякое энергетическое утверждение обязано идти в паре с именованным предикатом
интегрируемости --- здесь это не пожелание, а доказанное предупреждение.

\paragraph{От монолита к точечному балансу.}
Аналитическое ядро ветки --- энергетическое неравенство
$E(u(t_2)) + \int_{t_1}^{t_2} D(u(s))\,ds \le E(u(t_1))$, названное \code{TwoTimeEnergyInequality}.
Работа его не доказала, а \emph{сузила} до одного точечного входа \code{EnergyBalanceLaw},
\code{HasDerivAt}-тождества $dE/dt = -D(t)$ --- классического энергобаланса гладкого быстро
убывающего решения, в котором конвекция и давление не работают благодаря $\operatorname{div} u = 0$.
Глюа от узкого входа доказана.

\begin{theorem}[\code{twoTimeEnergyInequality\_of\_energyBalance}]\label{thm:twoTimeEnergyInequality_of_energyBalance}
Если выполнен точечный баланс $\code{EnergyBalanceLaw}\ \nu\ u$ и $s \mapsto D(u(s))$ интервально
интегрируемо на всех $[t_1,t_2]$, то выполнено $\code{TwoTimeEnergyInequality}\ \nu\ u$. \green{}
\end{theorem}

\emph{Идея доказательства.} Точечный баланс плюс интегрируемость дают двухвременное неравенство (на
деле равенство) через основную теорему анализа
(\code{intervalIntegral.integral\_eq\_sub\_of\_hasDerivAt}). Вход стал строго \'{у}же: не интегральное
неравенство по всем парам времён, а одно дифференциальное тождество.

\paragraph{Цепь до каскада.}
Пусть $\code{DissipativeStage}\ \nu\ u\ \delta\ t_1\ t_2$ отмечает временной интервал, на котором
накоплено диссипации не меньше $\delta$. Бюджетная машина запрещает бесконечную башню таких шагов.

\begin{theorem}[\code{ns\_no\_infinite\_dissipative\_cascade}]\label{thm:ns_no_infinite_dissipative_cascade}
Пусть $\delta > 0$ и выполнено $\code{TwoTimeEnergyInequality}\ \nu\ u$. Тогда не существует
последовательности времён $(t_k)_{k\in\N}$ с $\code{DissipativeStage}\ \nu\ u\ \delta\ t_k\ t_{k+1}$
для всех $k$. \green{} (условно на входе-неравенстве)
\end{theorem}

\emph{Идея доказательства.} Накопленная плата $\sum_k \delta$ превысила бы стартовую энергию
$E(u(t_0))$; прямое применение квантованной каскадной леммы
\code{no\_infinite\_uniform\_dissipative\_cascade}. Это ровно тот сертификат, которого требует
регулярность --- \emph{если} посылку кто-то предъявит. Слово «квантованной» несёт всю нагрузку, и
причина зафиксирована машинно.

\begin{theorem}[\code{real\_positive\_work\_not\_wellfounded}]\label{thm:real_positive_work_not_wellfounded}
Существует $a : \N \to \R$ такое, что $a_{n+1} < a_n$, $a_n > 0$ и $a_n - a_{n+1} > 0$ для всех $n$:
$\R$-последовательность со строго положительной работой на каждом шаге, но конечным суммарным спадом.
\green{}
\end{theorem}

\emph{Идея доказательства.} Банальный свидетель $a_n = 1/2^n$. Для $\N$-движка близнецов строго
положительная работа даёт well-foundedness даром; для $\R$ --- нет. Потому аналогия с близнецовым
движком лишь структурна: конечность каскада покупается равномерной нижней гранью $\delta > 0$, а не
положительностью как таковой. Откуда взять $\delta$ для реального решения --- аналитика, не
структура, и это остаётся открытым.

\paragraph{Интеграл взят.}
FTC-глюа поднимается до точного тождества, а само решение восстанавливается интегрированием.

\begin{theorem}[\code{energy\_identity\_of\_energyBalance}]\label{thm:energy_identity_of_energyBalance}
При $\code{EnergyBalanceLaw}\ \nu\ u$ и интервальной интегрируемости диссипации для всех $t_1,t_2$
выполнено
\begin{equation}\label{eq:ns-energy-identity}
E(u(t_2)) = E(u(t_1)) - \int_{t_1}^{t_2} D(u(s))\,ds,
\end{equation}
равенством; двухвременное неравенство
\cref{thm:twoTimeEnergyInequality_of_energyBalance} --- его огрубление. \green{}
\end{theorem}

Гипотетическая сингулярность за конечное время --- это \code{SingularCascade}: бесконечно много
$\delta$-диссипативных стадий, сжатых к одному конечному моменту $T$, --- тот же вечный двигатель,
что и всюду. Его невозможность зелена.

\begin{theorem}[\code{noSingularCascade\_of\_energyBalance}]\label{thm:noSingularCascade_of_energyBalance}
Точечный энергобаланс $dE/dt = -D$ вместе с интегрируемостью влечёт отсутствие точек разрыва
каскадного типа: $\code{NoSingularCascade}\ \nu\ u$. \green{}
\end{theorem}

\emph{Идея доказательства.} Прямое применение \cref{thm:ns_no_infinite_dissipative_cascade}: бюджет
запрещает бесконечный $\delta$-каскад, а сжатый к конечному времени --- тем более. Подчеркнём
честную оговорку: \code{NoSingularCascade} --- это \emph{не} $C^\infty$-регулярность, а отсутствие
бесконечной равномерно-квантованной диссипативной лестницы, то есть суррогат. Что оговорка
«равномерно» несущая, показывает кованый неравномерный каскад
(\code{cookedProfileCascade\_not\_uniform}, \green{}): его падения ускользают под всякий $\delta$, и
бюджет его не берёт.

\paragraph{На ℝ³: интеграл собран на боксе.}
Выше энергия жила абстрактным профилем, переданным гипотезой. На классе бокс-носимых решений баланс
уже не постулируется, а \emph{выводится}, задействуя теорему \code{mathlib} о дивергенции (в
box-форме \code{integral\_divergence\_of\_hasFDerivAt\_off\_countable}) и дифференцирование под
интегралом, через три тождества интегрирования по частям (давление и конвекция гаснут по
бездивергентности; вязкость даёт минус энстрофию).

\begin{theorem}[\code{noSingularCascade\_of\_r3Assembly}]\label{thm:noSingularCascade_of_r3Assembly}
Для бокс-носимого $C^2$-решения с силой $f = 0$ каскадная гладкость \code{NoSingularCascade}
выводится зелёно, без декрета первопричины \code{nsBoundary}. \green{}
\end{theorem}

\emph{Идея доказательства.} На боксе \code{energyBalance\_of\_boxSupported} выводит $dE/dt = -D$ из
производной под интегралом и трёх тождеств интегрирования по частям, после чего работает бюджетная
машина \cref{thm:noSingularCascade_of_energyBalance}; \code{\#print axioms} показывает стандартную
тройку, ни \code{sorry}, ни новой аксиомы. За пределами бокс-класса предельный переход
бокс $\to \R^3$ для лишь убывающих (не финитных) полей в \code{mathlib} отсутствует --- внешний
разрыв --- и \code{EnergyBalanceLaw} в общем случае остаётся открытым \red{}-входом.

\paragraph{Редукция Клэя: один недостающий критерий.}
Точная постановка Клэя-(A) --- данные Шварца, бездивергентность, существование гладкого глобального
решения с конечной энергией --- закодирована и машинно сведена к единственному открытому утверждению.

\begin{theorem}[\code{clayA\_of\_regularityTransfer\_and\_vorticityControl}]\label{thm:clayA_of_regularityTransfer_and_vorticityControl}
Известная теория переноса --- локальное сильное существование Като \cite{kato1984} вместе с
критерием продолжения Билла--Като--Майды \cite{bkm1984} --- и единственный критерий
\code{GlobalVorticityControl} влекут Клэй-(A). \green{}
\end{theorem}

\emph{Идея доказательства.} Чистый модус поненс: локальное гладкое решение существует; критерий не
даёт вихрю взорваться на каждом конечном промежутке; \cite{bkm1984} продлевает решение на всё время.
Теорема зелёная и \emph{не вакуумная} --- заключение обитаемо (\code{globalSmoothSolution\_zero}), а
оба входа реально потребляются. Смысл --- изоляция: всё, что осталось для (A), стянуто в одно
именованное открытое ядро \code{GlobalVorticityControl}, барьер суперкритичности, доказать который не
умеет никто. Саму теорию переноса (\code{RegularityTransfer}, \code{RegularityNecessity}) мы честно
называем \red{}-входом --- это реальные теоремы литературы \cite{kato1984,bkm1984}, не формализованные
в \code{mathlib}. Соблазнительно закрыть (A) через двигатель, как можно для Коллатца; у Навье--Стокса
этот мост доказуемо рвётся, и модуль фиксирует это машинно.

\begin{theorem}[\code{greenBudget\_strictly\_weaker\_than\_vorticityControl}]\label{thm:greenBudget_strictly_weaker_than_vorticityControl}
Зелёная бюджетная машина строго слабее \code{GlobalVorticityControl}: она не влечёт критерий. \green{}
\end{theorem}

\emph{Идея доказательства.} Суррогат \code{NoSingularCascade} убивает лишь равномерно-квантованный
каскад, а \code{cookedProfileCascade\_not\_uniform} предъявляет неравномерный каскад, который
ускользает. Настоящий взрыв --- континуальное явление; дискретный запрет двигателя его не достаёт.
Честность требует именовать открытое ядро аналитически (вихрь), а не прятать его в метафору. Эта
работа поэтому \emph{не приближает} решение Навье--Стокса ни на шаг математически: она даёт точную
формулировку, проверенный скелет редукции, точное имя единственной недостающей теоремы и барьер,
из-за которого та вне досягаемости. Навье--Стокс \emph{не} решён.

\paragraph{Диадическая модель: где двигатель работает и где взрывается.}
Чтобы проверить прочтение-двигатель на строгой дискретной жидкости, мы формализуем диадическую модель
Кац--Павлович \cite{katz-pavlovic2005}. (Ни одной модели жидкости прежде не было формализовано ни в
одном пруф-ассистенте; новизна формализационная, не математическая --- ничего здесь не решается.)
Сама стратегия --- добиться взрыва за конечное время в тщательно модифицированном уравнении жидкости
--- в духе усреднённого трёхмерного взрыва {Navier--Stokes} Тао~\cite{tao2016}.
Где диссипация равномерна --- бюджет побеждает; где каскад настоящий --- двигатель
\emph{реализуется}, и наивное «двигателя нет $\Rightarrow$ сингулярности нет» опровергнуто.
Ригорозное ядро --- факт об ОДУ.

\begin{theorem}[\code{superlinear\_blowup\_sq}]\label{thm:superlinear_blowup_sq}
Ни одна глобальная положительная $C^1$-функция не может удовлетворять $y'(t) \ge C\,y(t)^2$ с
$C > 0$: предположение о глобальном решении даёт \code{False}. \green{}
\end{theorem}

\begin{proof}
Положим $w(t) := 1/y(t) + C\,t$. Тогда $w'(t) = -y'/y^2 + C \le -C + C = 0$, значит $w$ не
возрастает, откуда $1/y = w - C\,t$ обязана стать отрицательной за конечное время --- противоречие с
$y > 0$. Супер-линейный каскад --- реализованный вечный двигатель: он существует ровно до момента
взрыва.
\end{proof}

\begin{theorem}[\code{sublinear\_collapse\_extinction}]\label{thm:sublinear_collapse_extinction}
\green{}\ Двойственно: нет глобальной положительной $C^1$-функции с $r'(t) \le -C\,r(t)^{\alpha}$ при
$C>0$, $0 \le \alpha < 1$; предположение о глобальном решении даёт \code{False}. Коллапс к $r=0$
наступает за конечное время, не позднее $r(0)^{1-\alpha}/((1-\alpha)C)$.
\leantag{sublinear\_collapse\_extinction}
\end{theorem}

\begin{proof}
Положим $v(t) := r(t)^{1-\alpha} + (1-\alpha)C\,t$. Тогда
$v'(t) = (1-\alpha)r^{-\alpha}r' + (1-\alpha)C \le 0$ (из $r'\le -C\,r^{\alpha}$ и
$r^{\alpha}r^{-\alpha}=1$), значит $v$ не возрастает, откуда $r^{1-\alpha} = v - (1-\alpha)C\,t$
обязана стать отрицательной за конечное время --- противоречие с $r>0$.
\end{proof}

\emph{Почему это точный дуал.} Взрыв и коллапс --- один драйв $u' = \pm C\,u^{\beta}$ по разные
стороны критического показателя $\beta=1$: сверхлинейный ($\beta>1$) уходит к $\infty$, сублинейный
($\beta<1$) --- к $0$, при $\beta=1$ конечно-временной сингулярности нет. Одна подстановка
$w := u^{1-\beta}$ линеаризует оба, давая тот же антитонный помощник и тот же шаг о среднем. Новизна
формализационная, не математическая: факт вымирания классичен, а доказательство --- зеркало взрыва
выше. Вывод драйва $r'\le -C\,r^{\alpha}$ из физического закона само-гравитации (Эйлер--Пуассон,
автомодельный коллапс Ларсона--Пенстона) --- именованный аналитический вход за пределами
\code{mathlib}, \emph{здесь не сделан}; прочтения «коллапс / сверхновая / Большой взрыв» --- метафора,
не физика. Дуальное приложение-модель, а не маска-миллениум.

\begin{theorem}[\code{dyadic\_blowup}]\label{thm:dyadic_blowup}
Модель Кац--Павлович \code{DyadicSolution} с уравнениями
$a_n' = \lambda^n a_{n-1}^2 - \lambda^{n+1} a_n a_{n+1} - d_n a_n$ глобально пуста: взрыв неизбежен.
\green{}
\end{theorem}

\emph{Идея доказательства.} Нелинейная передача телескопируется (энергия сохраняется связью), а взрыв
следует из \cref{thm:superlinear_blowup_sq}, применённой к ведущей моде. Супер-линейный драйв
$y' \ge C\,y^2$ больше не постулируется: для целого класса решений он выведен из
$\lambda^n$-связей.

\begin{theorem}[\code{frontDrive\_of\_invariant}]\label{thm:frontDrive_of_invariant}
Пусть у решения Кац--Павлович одна фронт-оболочка снизу зажата двумя соседями --- инвариант
\code{FrontDomination}: $\rho\, a_{J+1} \le a_J$, $a_{J+2} \le \kappa\, a_{J+1}$, $m \le a_{J+1}$.
Тогда драйв $C\,y^2 \le y'$ выведен прямо из $\lambda^n$-связей, с
$C = \lambda^{J+1}\rho^2 - \lambda^{J+2}\kappa - d_{J+1}/m$. \green{}
\end{theorem}

\emph{Идея доказательства.} Три оценки подставляются в правую часть оболочки $J+1$: приток
$\lambda^{J+1} a_J^2$ рождает $y^2$, а отток и диссипация вычитаются под контролем. Не-вакуумность
подтверждена точным автомодельным профилем $a_n(t) = \lambda^{-n}/((\lambda^2-1)(T-t))$,
удовлетворяющим инварианту. Один вход остаётся открытым: сохранность \code{FrontDomination} для
\emph{бесконечно} многих мод на всём времени жизни --- движущийся фронт Кац--Павлович,
исследовательский фронтир, здесь не доказанный.

Исток каскада несёт единственный намеренный жёлтый слой этой ветки. У нижней оболочки $n = 0$ приток
нулевой: она умеет только отдавать, сама себя не запустит, и её исток нельзя породить изнутри связей.
Автомодельному решению нужен внешний накачивающий член, и эта поставка берётся из первопричины.

\begin{theorem}[\code{dyadicBlowup\_is\_firstCauseManifestation}]\label{thm:dyadicBlowup_is_firstCauseManifestation}
Поставка на масштабе $n = 0$ берётся из границы \code{nsBoundary} аксиомы \axiom{}: диадический взрыв
присоединяется к манифестациям первопричины через свой исток. \yellow{}~(\textsc{axiom-tainted})
\end{theorem}

\emph{Идея доказательства.} Две стены: зелёная стена невозможности (исток неказуем изнутри) плюс
поставка, декретированная извне. Зелёный вывод драйва самодостаточен и от этого слоя не зависит ---
убери первопричину, и математика взрыва останется целой, исчезнет лишь ответ на вопрос «кто зажёг
насос $n=0$», --- поэтому таинт растёт ровно на две декларации и ни на одну больше.

\paragraph{Не Большой взрыв, а одно из его лиц.}
Исток легко перечитать космологически чересчур буквально; мы этого не делаем. В космологическом
переводе Большим взрывом служит сама первопричина, событие $0\to1$; Навье--Стокс же --- лишь одна из
трёх её границ (\code{step00FirstCause\_iff\_causalClosure}: первопричина эквивалентна конъюнкции
узла близнецов, закона Римана и гейта НС). Жидкостный взрыв не есть начало мира --- его исток $n=0$,
не умея завестись сам, лишь занимает искру у того же единственного начала: одно из трёх лиц одной
сингулярности, а не сама сингулярность. Как всюду в программе, космология здесь --- перевод теорем,
а не претензия на физику.

\paragraph{Честный статус.} Прочтение-двигатель не оправдано и не отброшено: оно
\emph{картировано}. Где диссипация равномерна, бюджет запрещает вечный каскад; где каскад настоящий,
двигатель реализуется и жидкость взрывается в модели. Новой математики для проблемы тысячелетия не
произведено (это тень барьера суперкритичности Тао \cite{clay-ns}), лишь первая в мире формализация
дискретной модели жидкости из \emph{известной} математики, точный скелет редукции и машинная карта
того, где прочтение справедливо и где оно ломается. Навье--Стокс \emph{не} решён.
